% ============================================================
% figure_caption_oral_absorption.tex
% ============================================================
% Publication-quality LaTeX caption for the 4-panel oral
% absorption figure (plot_oral_absorption_model.py).
%
% USAGE IN YOUR MANUSCRIPT
% ─────────────────────────
% 1. Copy the \begin{figure}...\end{figure} block into your
%    manuscript at the correct position.
% 2. Add the three BibTeX entries at the bottom of this file
%    to your .bib file.
% 3. Adjust the \includegraphics path to match your figure
%    filename/location.
% 4. Compile with:
%       pdflatex manuscript
%       bibtex   manuscript
%       pdflatex manuscript  (×2)
%
% NOTES ON MATH INSIDE \caption{}
% ─────────────────────────────────
% LaTeX does not support \begin{equation*} inside \caption{}.
% All equations here are written as inline math ($...$) which
% compiles cleanly and is standard practice for figure captions.
% ============================================================

\begin{figure}[htbp]
    \centering
    \includegraphics[width=\textwidth]{figures/oral_absorption_model_TID}
    \caption{%
        \textbf{Oral absorption model for furosemide:
        piecewise-constant input rate approximation and
        sensitivity analysis.}
        %
        \textbf{(A)} Amount of drug remaining in the gut compartment
        as a function of time after a single oral dose, described by
        the first-order absorption equation
        $A_{\mathrm{gut}}(t) = F \cdot D \cdot e^{-k_a t}$,
        where $F = 0.60$ (oral bioavailability \cite{FDA_furosemide_label}),
        $D = 240$\,mg (total daily dose, TID regimen: $3 \times 80$\,mg),
        and $k_a = 1.5$\,h$^{-1}$ \cite{Hammarlund1984}.
        The shaded area denotes the total absorbed amount over the
        6-hour dosing window; the highlighted segment illustrates a
        single piecewise interval $j$ with absorbed amount
        $A_{\mathrm{abs},j}$.
        %
        \textbf{(B)} Absorption rate as a function of time.
        The continuous curve represents the analytical rate
        $\dot{A}_{\mathrm{abs}}(t) = k_a \cdot A_{\mathrm{gut}}(t)$;
        bars represent the piecewise-constant infusion rate
        approximation $IR_j = A_{\mathrm{abs},j}/\Delta t_j$
        across $N = 20$ equal time segments, used as the drug input
        to the ODE system.
        Total absorbed dose: $\sum A_{\mathrm{abs}} = 143.98$\,mg
        (0.012\% error relative to $F \cdot D$).
        %
        \textbf{(C)} Summary of the governing equations.
        The gut compartment obeys
        $dA_{\mathrm{gut}}/dt = -k_a A_{\mathrm{gut}}$,
        with analytical solution
        $A_{\mathrm{gut}}(t) = F \cdot D \cdot e^{-k_a t}$.
        The absorbed amount in segment $j$ is
        $A_{\mathrm{abs},j} = F \cdot D
        \left[e^{-k_a t_j} - e^{-k_a t_{j+1}}\right]$,
        and the corresponding constant infusion rate delivered to the
        central compartment is $IR_j = A_{\mathrm{abs},j}/\Delta t_j$.
        %
        \textbf{(D)} Sensitivity of the gut absorption profile to the
        absorption rate constant $k_a$.
        Three values are shown: $k_a = 1.0$, $1.5$ (central
        estimate \cite{Hammarlund1984}), and $2.0$\,h$^{-1}$.
        Vertical dashed lines indicate the time of peak plasma
        concentration ($T_{\max}$), determined numerically from the
        two-compartment oral model by solving $dC_{\mathrm{cen}}/dt = 0$,
        using the macro rate constants $\alpha = 0.631$\,h$^{-1}$ and
        $\beta = 0.374$\,h$^{-1}$ derived from the CKD-5 model parameters
        ($k_{10} = CL_{\mathrm{furo}}/V_{\mathrm{cen}} = 0.472$\,h$^{-1}$,
        $k_{12} = 0.033$\,h$^{-1}$, $k_{21} = 0.50$\,h$^{-1}$,
        $CL_{\mathrm{furo}} = 2.36$\,L\,h$^{-1}$, $V_{\mathrm{cen}} = 5.0$\,L).
        Predicted $T_{\max}$ values of 84, 66, and 56\,min for
        $k_a = 1.0$, $1.5$, and $2.0$\,h$^{-1}$, respectively, are
        consistent with the reported oral furosemide tablet
        $T_{\max}$ of approximately 87\,min in healthy fasted
        subjects \cite{FDA_furosemide_label}.
    }
    \label{fig:oral_absorption_model}
\end{figure}

% ============================================================
% BibTeX entries — add these to your .bib file
% ============================================================

% FDA label — primary source for F=0.60 and Tmax~87 min
@manual{FDA_furosemide_label,
  title        = {{LASIX\textsuperscript{\textregistered}
                   (furosemide) Tablets 20, 40, and 80\,mg:
                   Prescribing Information}},
  organization = {{U.S. Food and Drug Administration}},
  note         = {NDA 016273, Revised 2012.
                  Available: \url{https://www.accessdata.fda.gov/
                  drugsatfda_docs/label/2012/016273s066lbl.pdf}},
  year         = {2012},
}

% Hammarlund 1984 — source for ka and absorption-rate-limited kinetics
@article{Hammarlund1984,
  author  = {Hammarlund, Margareta M. and
             Paalzow, Lennart K. and
             Odlind, Birger},
  title   = {Pharmacokinetics of furosemide in man after intravenous
             and oral administration. {Application} of moment analysis},
  journal = {European Journal of Clinical Pharmacology},
  year    = {1984},
  volume  = {26},
  number  = {2},
  pages   = {197--207},
  doi     = {10.1007/BF00630286},
  pmid    = {6723758},
}

% Rowland & Tozer — textbook source for first-order absorption equation form
@book{Rowland2011,
  author    = {Rowland, Malcolm and Tozer, Thomas N.},
  title     = {Clinical Pharmacokinetics and Pharmacodynamics:
               Concepts and Applications},
  edition   = {4th},
  publisher = {Lippincott Williams \& Wilkins},
  address   = {Philadelphia},
  year      = {2011},
  note      = {Chapter~4: Multiple-Dose Regimens},
}
